\paragraph{Reading an added vertex.}
For the actual full-two-torsion family, all primes are split,
$\delta\in\mathbb Z$, and $n_U=\delta+|U|$.  Put $F_U=\theta H_U$
and let $q\notin S\cup U$ be a fresh prime, where $S$ includes
every late sieve and the current moving theta prime.
Write $u=\prod_{p\in U}p$ and let $\mathcal A(D)$ denote the
normalized valuation symbol on the tested range: it is one exactly when
$v_2a(D)=\delta+\omega(D)$ on the filter. For an odd source, that
range has a fixed weight bound chosen before the valid actual stage,
as specified before (F.7). No finite-stage value at larger weights is
asserted; the compatible limit tables used later are a separate object.
Expanding the operators, every term containing a $V_p$ has zero
coefficient at $q$, and no unary index $Jm^2$ with $J\subset U$
equals $q$.  Therefore
\[
 B_U[q]=2^{-n_U}a(uq),\qquad
 H_U[q]=2^{-n_U-1}a(uq).
\]
At coefficient one the diamond term in the good Hecke formula
vanishes.  Since $F_U\equiv H_U\pmod{2\mathcal O_{\rm coef}}$,
\[
 \overline{(t(\operatorname{Fr}_q)F_U)[1]}
   =\overline{H_U[q]}
   =\overline{\,2^{-\delta-|U|-1}a(uq)\,}
   =\mathcal A(uq)\in\mathbb F_2.
\]
Here residues lie in the coefficient residue field, with
$\mathbb F_2$ its prime subfield.  Every permitted phase is a
two-power root of unity and reduces to one.  For an odd source,
fix the finite operator budget $|U|\le W$ and $R\ge2$ first;
the stage choice in (O.9) may be taken with
$M_\theta\ge\max(0,\delta+W)+R$.
Thus $M_\theta-n_U-1\ge1$, so the normalized finite-stage error
is still divisible by the integer two.  The same reading holds
with value zero off the fixed or moving coefficient support.

The diamond correction above is a sum of integral multiples of
$\theta\Theta_J$; its fresh coefficient $q$ is zero after reduction.
It introduces neither an extra phase nor a denominator in this
reading.  All readings and subsequent omissions use this same
$g,\eta,\psi,\delta$ and actual stage; only the set $U$ changes.

For any fixed Fourier index $n$, Chebotarev approximation lets one
test a finite product of trace operators using distinct fresh primes
$q_1,\ldots,q_a$ avoiding $n$ and the level. The good Hecke formula
then reads the coefficient at $nq_1\cdots q_a$ of the original vector.
Different indices may use different fresh primes; the conclusion is
an identity of coefficientwise reduced expansions.

